\lecture{6}
\title{Uncertainty, Regularizers, and Features}

\begin{document}

\subsection{Aleatoric vs. Epistemic Uncertainty}

There are two types of uncertainty
that we would like to conceptually differentiate here.

\textbf{Aleatoric Uncertainty.}
This refers to uncertainty caused by
the intrinsic randomness in how the response $y^{(n)}$
relates to the feature vector $x^{(n)}$.
It's the uncertainty in the PDF $p(y \mid x, \theta)$.

\textbf{Epistemic Uncertainty.}
This refers to uncertainty caused by
an incomplete model of the data being perceived.
It's the uncertainty in the posterior $p(\theta \mid X, Y)$
caused by an insufficient amount of data.

Reading in more data will generally lessen the
epistemic uncertainty, but more data won't be able to help
with aleatoric uncertainty. As
\href{/6.790/lectures/5/#bayesian-linear-regression-performance}{Lecture 5}
showed, a misspecified model will produce great aleatoric uncertainty,
yet the posterior overconfidently presents itself with
very high epistemic certainty as it receives more data.

\begin{remark}
    This overconfidence isn't unique to the Bayesian setting, of course;
    in the frequentist setting, standard errors and
    CI-widths also generally decrease as more data is read,
    even if the aleatoric uncertainty remains high.
\end{remark}

\subsection{MAP for Bayesian Linear Regression}

We'll make some simplifying assumptions about our prior:
assume $\theta \sim \mathcal{N}(0, \sigma_0^2 \cdot I_D)$,
where $D$ is the feature dimension.
Then the posterior looks like the following; compare with the proof in
\href{/6.790/lectures/5/#bayesian-linear-regression-theory}{Lecture 5},
perhaps.
\[ p(\theta \mid X, Y)
\propto \underbrace{\mathrm{exp}\left(
    -\frac{\|Y - X\theta\|^2}{2\sigma^2}
\right)}_{
    \propto ~ L_N(\theta)
} \cdot \underbrace{\mathrm{exp}\left(
    -\frac{\|\theta\|^2}{2\sigma_0^2}
\right)}_{\propto ~ p(\theta)}. \]
Then the maximum a posteriori (MAP)
satisfies the following:

\begin{theorem}{Bayesian MAP}
    The maximum a posteriori for Bayesian linear regression,
    under the prior $\theta \sim \mathcal{N}(0, \sigma_0^2 I_D)$, is:
    \[ \hat{\theta}_{\mathrm{MAP}}
    = \argmin_{\theta} \left[ \|X\theta - Y\|^2
    + \frac{\sigma^2}{\sigma_0^2}\|\theta\|^2 \right]
    = \left(\frac{\sigma^2}{\sigma_0^2}I_D + X^{\top}X\right)^{-1}
    X^{\top}Y. \]
\end{theorem}

\begin{proof}
    The first equality is trivial,
    so it remains to show the RHS expression for $\hat{\theta}_{\mathrm{MAP}}$
    minimizes the middle expression.
    A shortcut is to just reuse the results we derived from
    \href{/6.790/lectures/5/#bayesian-linear-regression-theory}{Lecture 5},
    but for $(\mu_0, \Sigma_0) = (0, \sigma_0^2 I_D)$.
    \[ \Sigma_N^{-1}
    = \frac{1}{\sigma_0^2} I_D + \frac{1}{\sigma^2} X^{\top}X
    ~~~ \text{ and } ~~~
    \Sigma_N^{-1}\mu_N
    = \frac{1}{\sigma^2}X^{\top}Y. \]
    The posterior distribution is a Gaussian $\mathcal{N}(\mu_N, \Sigma_N)$,
    so it is maximized at $\hat{\theta}_{\mathrm{MAP}} = \mu_N$.
    The result falls right out from there. ~ $\blacksquare$
\end{proof}

The quantity being minimized in the theorem statement above
is referred to as the \textbf{ridge regression objective}.

It is the sum of our ordinary squared-loss
and a new ``ridge regression term'',
the latter of which penalizes point estimates $\hat{\theta}_{\mathrm{MAP}}$
that have needlessly large coefficients---large
coefficients in $\hat{\theta}_{\mathrm{MAP}}$
often suggest overfitting.

We'll oftentimes denote $\lambda := \frac{\sigma^2}{\sigma_0^2}$
in the ridge regression objective and
expression for $\hat{\theta}_{\mathrm{MAP}}$.
Note that, while we can choose the value of $\sigma_0$ in our prior,
we usually have no control or knowledge regarding $\sigma$.
So the value of $\lambda$ can alternatively be interpreted
as a parameter of our choosing that indicates
how strongly we want to discourage overfitting.

\subsection{Features}

Let's stop restricting ourselves to linear models.
The idea is to create our own feature vectors
whose entries are nonlinear functions
of the entries of the feature vectors $\{x^{(n)}\}_{n = 1}^N$.

\textbf{Definition.}
We define our own \textbf{features} as
$\phi: \mathbb{R}^D \to \mathbb{R}^P$
and write our regression hypothesis as
$h(x) = \theta^{\top} \phi(x)$.

\textbf{Example.}
In the case of $D = 1$,
we might take $\phi(x) := \begin{bmatrix} 1 & x & x^2 \end{bmatrix}^{\top}$,
which yields the hypothesis
$h(x) = \theta_0 + \theta_1x + \theta_2x^2$.

\textbf{Example.}
In the case of $D = 2$,
we might take $\phi(x) := \begin{bmatrix}
1 & x_1 & x_2 & x_1^2 & x_2^2 & x_1x_2
\end{bmatrix}^{\top}$,
and the set of all possible hypotheses
becomes the set of all two-variable quadratic polynomials.

Here's how regression performs
when we get nonlinear features involved.

\begin{center}
    \includegraphics[width=0.75\textwidth]{lecture-06/feature-examples}
\end{center}

Notice how setting the feature dimension $P$ too large
results in overfitting, to the point where
our degree-$9$ regression is just performing
Lagrange interpolation.
To reduce overfitting, let's bring back a trick
from just a few paragraphs ago: ridge regression.

\begin{center}
    \includegraphics[width=0.3\textwidth]{lecture-06/ridge-example}
\end{center}

And so degree-$9$ regression works just fine
with a ridge regression term using $\lambda = e^{-18}$.

\subsection{Kernels: Unlimited Features}

There's another reason why too many features is bad:
if we let $\Phi \in \mathbb{R}^{N \times P}$
(with rows $\phi(x^{(n)})^{\top}$)
be our stand-in for $X \in \mathbb{R}^{N \times D}$,
then the ridge solution looks like:
\[ \hat{\theta}_{\mathrm{MAP}}
= \left(
    \lambda I_P + \Phi^{\top}\Phi
\right)^{-1} \Phi^{\top}Y. \]
Computing this would require a $P \times P$ matrix inversion,
which is far too computationally intensive for $P \gg D$.
So it seems like there is an upper bound on
the number of features $P$ we can reasonably have.

Or maybe there isn't!  We claim that there is an alternative way
to approach our computation of ridge regressions
that allows for $P$ to be arbitrarily large
at no computational cost;
the rest of this section
is dedicated to showing how.

First, we claim the following theorem is true;
the proof is unenlightening algebraic manipulation, so we omit it.

\begin{theorem}{Bayesian MAP, Rewritten}
    The ridge solution $\hat{\theta}_{\mathrm{MAP}}$ may be expressed as:
    \[ \hat{\theta}_{\mathrm{MAP}}
    = \Phi^{\top}\left(
        \lambda I_N + \Phi\Phi^{\top}
    \right)^{-1}Y. \]
    Note, importantly, that this implies that the
    ridge regression may be expressed as:
    \[ h(x) = \hat{\theta}^{\top}_{\mathrm{MAP}} \phi(x)
    = Y^{\top}\left(
        \lambda I_N + \Phi\Phi^{\top}
    \right)^{-1} \Phi \phi(x). \]
\end{theorem}

In this form, it turns out that there's
a great amount of simplification to be made
if we introduce the following:

\textbf{Definition.}
The \textbf{kernel} is the function
$k \colon \mathcal{X} \times \mathcal{X} \to \mathbb{R}$ (where $\mathcal{X}$
is the input space) given by
$k(x^{(i)}, x^{(j)}) := \phi^{\top}(x^{(i)}) \phi(x^{(j)})$.
\begin{itemize}
    \item Given data $\{x^{(n)}\}_{n = 1}^N$,
    we may define a map $k: \mathcal{X} \to \mathbb{R}^N$
    via $k(x)_i := k(x^{(i)}, x)$.
    (Yes, overloaded notation.)
    \item Given data $\{x^{(n)}\}_{n = 1}^N$,
    we may define a matrix $K \in \mathbb{R}^{N \times N}$
    given by $K_{i, j} := k(x^{(i)}, x^{(j)})$.
\end{itemize}

We can now again rewrite the ridge regression formula
using the language of the kernel.

\begin{theorem}{Bayesian MAP, w/ Kernel}
    The ridge regression may be expressed as:
    \[ h(x) = Y^{\top}\left(\lambda I_N + K\right)^{-1} k(x). \]
\end{theorem}

Here's the upshot of all of this:
the expression above only deals with
matrices in $\mathbb{R}^{N \times N}$
and vectors in $\mathbb{R}^{N}$;
there's no dependency on $P$ anywhere!
So as long as we can efficiently compute
the kernel $k \colon \mathcal{X} \times \mathcal{X} \to \mathbb{R}$,
we can set our number of features $P$ as large as we'd like.

\end{document}